% R2_Decomposition_Theorem.tex
% monogate Cost Theory — T38: Cost Decomposition Theorem
% Formal proof of Cost(E) = NaiveCost(E) - SharingDiscount(E) - PatternBonus(E)
% Classification: THEOREM (T38)
% Author: Arturo R. Almaguer, 2026-04-20
% Encoding: UTF-8

\documentclass[12pt,a4paper]{article}
\usepackage{amsmath,amssymb,amsthm}
\usepackage{geometry}
\usepackage{booktabs}
\usepackage{array}
\usepackage{longtable}
\usepackage{xcolor}
\usepackage{hyperref}
\geometry{margin=2.5cm}

% ── Theorem environments ────────────────────────────────────────────────────
\newtheorem{theorem}{Theorem}[section]
\newtheorem{lemma}[theorem]{Lemma}
\newtheorem{corollary}[theorem]{Corollary}
\newtheorem{proposition}[theorem]{Proposition}
\theoremstyle{definition}
\newtheorem{definition}[theorem]{Definition}
\newtheorem{conjecture}[theorem]{Conjecture}
\newtheorem{remark}[theorem]{Remark}
\newtheorem{example}[theorem]{Example}
\newtheorem{observation}[theorem]{Observation}

% ── Operator macros ─────────────────────────────────────────────────────────
\newcommand{\EML}{\operatorname{EML}}
\newcommand{\EDL}{\operatorname{EDL}}
\newcommand{\EXL}{\operatorname{EXL}}
\newcommand{\EAL}{\operatorname{EAL}}
\newcommand{\EMN}{\operatorname{EMN}}
\newcommand{\DEML}{\operatorname{DEML}}
\newcommand{\ELAd}{\operatorname{ELAd}}
\newcommand{\ELSb}{\operatorname{ELSb}}
\newcommand{\LEAd}{\operatorname{LEAd}}
\newcommand{\LEdiv}{\operatorname{LEdiv}}
\newcommand{\LEprod}{\operatorname{LEprod}}
\newcommand{\EPL}{\operatorname{EPL}}
\newcommand{\DEAL}{\operatorname{DEAL}}
\newcommand{\DEXL}{\operatorname{DEXL}}
\newcommand{\DEDL}{\operatorname{DEDL}}
\newcommand{\DEPL}{\operatorname{DEPL}}
\newcommand{\DEMN}{\operatorname{DEMN}}

% ── Cost macros ──────────────────────────────────────────────────────────────
\newcommand{\Cost}{\operatorname{Cost}}
\newcommand{\NC}{\operatorname{NaiveCost}}
\newcommand{\SD}{\operatorname{SharingDiscount}}
\newcommand{\PB}{\operatorname{PatternBonus}}
\newcommand{\bonus}{\operatorname{bonus}}
\newcommand{\Fam}{\mathcal{F}}

% ── DAG / tree macros ────────────────────────────────────────────────────────
\newcommand{\Gmin}{G_{\min}(E)}
\newcommand{\Tnaive}{T_{\mathrm{naive}}(E)}
\newcommand{\Tunroll}{T_{\mathrm{unroll}}(E)}

% ── Column types ─────────────────────────────────────────────────────────────
\newcolumntype{L}[1]{>{\raggedright\arraybackslash}p{#1}}
\newcolumntype{C}[1]{>{\centering\arraybackslash}p{#1}}

\title{T38: The Cost Decomposition Theorem\\[6pt]
       \large
       $\Cost(E) = \NC(E) - \SD(E) - \PB(E)$\\[4pt]
       Formal Definitions, Proof, Uniqueness, and Worked Examples}
\author{Arturo R.\ Almaguer\\
\small monogate research \quad \texttt{almaguer1986@gmail.com}}
\date{April 2026}

% ════════════════════════════════════════════════════════════════════════════
\begin{document}
\maketitle

% ── Abstract ────────────────────────────────────────────────────────────────
\begin{abstract}
We give a self-contained formal treatment of Theorem T38, the central result
of the monogate SuperBEST cost theory.
T38 states that for every arithmetic expression~$E$ computable by a finite EML
operator tree over the 16-operator family~$\mathcal{F}_{16}$,
\[
  \Cost(E) \;=\; \NC(E) \;-\; \SD(E) \;-\; \PB(E),
\]
where the \emph{na\"{i}ve cost} $\NC(E)$ is the sum of per-operation SuperBEST
v3 unit costs assuming no sharing and no compound patterns, the
\emph{sharing discount} $\SD(E) \geq 0$ accounts for constant folding and
shared sub-expression elimination, and the \emph{pattern bonus} $\PB(E) \geq 0$
accounts for compound-operator compression.
Section~1 gives formal definitions of all three components.
Section~2 proves T38 by a four-step constructive argument and establishes that
the three terms are non-negative and exhaust all possible reductions.
Section~3 proves uniqueness of the decomposition for a fixed pattern set
$\mathcal{P}$ and canonical sharing strategy, and characterises how
the decomposition changes when $\mathcal{P}$ is extended.
Section~4 works three scientific examples in full detail: the Arrhenius rate
law (Class~A), the pH formula (Class~C), and the Maxwell-Boltzmann speed
distribution (Class~D).
\end{abstract}

\tableofcontents
\bigskip

% ════════════════════════════════════════════════════════════════════════════
\section{Formal Definitions}
\label{sec:defs}
% ════════════════════════════════════════════════════════════════════════════

% ── 1.1 SuperBEST v3 unit costs ─────────────────────────────────────────────
\subsection{SuperBEST v3 Unit Costs}

SuperBEST v3 assigns each primitive arithmetic operation a minimum node count
in $\mathcal{F}_{16}$.  These counts are optimal: no EML tree for the given
function uses fewer nodes.

\begin{table}[h]
\centering
\caption{SuperBEST v3 unit costs used throughout this paper.}
\label{tab:sb3}
\smallskip
\begin{tabular}{@{}llcc@{}}
\toprule
\textbf{Operation} & \textbf{Symbol} &
  \textbf{Cost $c_{\mathrm{op}}$} & \textbf{Notes} \\
\midrule
Exponential $e^x$            & \texttt{exp}   & 1  & $\EML(x,1)$ \\
Natural log $\ln x$          & \texttt{ln}    & 1  & $\LEdiv(x,1)$ \\
Division $x/y$               & \texttt{div}   & 1  & $\LEdiv(x,y)$ \\
Negation $-x$                & \texttt{neg}   & 2  & \\
Reciprocal $1/x$             & \texttt{recip} & 2  & \\
Multiplication $x\cdot y$    & \texttt{mul}   & 2  & \\
Subtraction $x - y$          & \texttt{sub}   & 2  & \\
Exponentiation $x^n$         & \texttt{pow}   & 3  & \\
Addition $x+y$ (all real $x,y$) & \texttt{add} & 2  & ADD-T1 (SuperBEST v5); positive/general split eliminated \\
\bottomrule
\end{tabular}
\end{table}

% ── 1.2 Na\"ive cost ─────────────────────────────────────────────────────────
\subsection{Na\"{i}ve Cost}

\begin{definition}[Na\"{i}ve Cost]
\label{def:NC}
For any arithmetic expression~$E$, the \emph{na\"{i}ve cost} is
\begin{equation}
  \NC(E) \;=\;
    1 \cdot n_{\exp}
  + 1 \cdot n_{\ln}
  + 1 \cdot n_{\mathrm{div}}
  + 2 \cdot n_{\mathrm{neg}}
  + 2 \cdot n_{\mathrm{rec}}
  + 2 \cdot n_{\mathrm{mul}}
  + 2 \cdot n_{\mathrm{sub}}
  + 3 \cdot n_{\mathrm{pow}}
  + 3 \cdot n_{\mathrm{add}},
\label{eq:NC}
\end{equation}
where $n_{\mathrm{op}}$ counts the occurrences of each primitive operation in
the expression \emph{tree} (with all shared sub-expressions unrolled and all
compound patterns replaced by their primitive expansions).
Numeric constants and free variables contribute~$0$.
\end{definition}

\begin{remark}[Interpretation]
$\NC(E)$ is the cost of the straightforward tree implementation: build one
independent sub-tree per operation occurrence, substituting each primitive with
its SuperBEST v3 optimal sub-tree, and make no attempt to merge equal
sub-expressions or to recognise compound patterns.
\end{remark}

% ── 1.3 DAG and unrolled tree ────────────────────────────────────────────────
\subsection{Minimal DAG and Unrolled Tree}

\begin{definition}[Minimal DAG]
\label{def:dag}
Let $\Gmin$ denote the \emph{minimal directed acyclic graph} (DAG) computing
$E$: the unique (up to isomorphism) DAG that
\begin{enumerate}
  \item correctly computes $E$ for all admissible inputs,
  \item uses SuperBEST v3 sub-trees for each primitive operation, and
  \item has the minimum number of nodes subject to (1) and (2).
\end{enumerate}
Minimality is achieved by sharing all identical sub-expressions and applying
all available compound-pattern substitutions.
\end{definition}

\begin{definition}[Unrolled Tree]
\label{def:unroll}
The \emph{unrolled tree} $\Tunroll$ is obtained from $\Gmin$ by replacing
every shared edge (i.e.\ every edge whose target node has in-degree $\geq 2$)
with an independent copy of the sub-DAG rooted at that node.  The result is a
tree (each node has in-degree $\leq 1$).
\end{definition}

\begin{remark}
$\Tunroll$ is, in general, larger than $\Tnaive$ because compound patterns
present in $\Gmin$ are still encoded as single nodes in $\Tunroll$; they are
expanded to their primitive equivalents only in $\Tnaive$.
\end{remark}

% ── 1.4 Sharing Discount ─────────────────────────────────────────────────────
\subsection{Sharing Discount}

\begin{definition}[Sharing Discount]
\label{def:SD}
Let $G = \Gmin$ and let $T = \Tunroll$ be the tree obtained by unrolling all
shared edges in~$G$.  Define:
\[
  \SD(E) \;=\; \NC(T) \;-\; \NC(G),
\]
where $\NC(\,\cdot\,)$ applied to a DAG or tree means the sum of
SuperBEST v3 node weights over all nodes (with compound-pattern nodes counted
at their primitive expansion cost in both cases, to make the subtraction
well-defined before pattern bonuses are separated out).

Equivalently, if node~$v$ in~$G$ is shared at $k_v \geq 2$ sites and has
sub-DAG cost $\Cost(v)$, then
\[
  \SD(E) \;=\; \sum_{\substack{v \in G \\ k_v \geq 2}} (k_v - 1)\cdot \Cost(v),
\]
where the sum is over all nodes with in-degree $\geq 2$ in~$G$.
\end{definition}

\begin{proposition}[$\SD(E) \geq 0$]
\label{prop:SD-nonneg}
For any expression $E$, $\SD(E) \geq 0$.
\end{proposition}

\begin{proof}
Every term $(k_v - 1)\cdot \Cost(v)$ in the sum satisfies $k_v \geq 2$ (so
$k_v - 1 \geq 1$) and $\Cost(v) \geq 1$ (since $v$ is a non-trivial node).
Hence every summand is non-negative, and the sum is non-negative.
\end{proof}

% ── 1.5 Pattern Bonus ────────────────────────────────────────────────────────
\subsection{Pattern Bonus}

\begin{definition}[Compound Pattern Set]
\label{def:patterns}
The \emph{compound pattern set} $\mathcal{P}$ for SuperBEST v3 is the finite
collection of pairs $(p_i, k_i^*)$ where $p_i \in \mathcal{F}_{16}$ is a
compound operator and $k_i^* = 1$ is its node count.  For each pattern $p_i$,
the \emph{na\"{i}ve cost} $k_i^{\mathrm{naive}}$ is the cost of implementing
$p_i$ using primitive operations without compression.

The standard patterns and their bonuses are:
\begin{center}
\begin{tabular}{@{}lcll@{}}
\toprule
\textbf{Pattern} & \textbf{Formula} & \textbf{Naive cost} &
  \textbf{Bonus} $= k_i^{\mathrm{naive}} - 1$ \\
\midrule
$\EML(x,y)$   & $e^x - \ln y$      & $1+1+2 = 4$   & 3 \\
$\EXL(x,y)$   & $e^x \cdot \ln y$  & $1+1+2 = 4$   & 3 \\
$\EAL(x,y)$   & $e^x + \ln y$      & $1+1+3 = 5$   & 4 \\
$\DEML(x,y)$  & $e^{-x} - \ln y$   & $2+1+1+2 = 6$ & 5 \\
$\LEAd(x,y)$  & $\ln(e^x + y)$     & $1+3+1 = 5$   & 4 \\
$\ELAd(x,y)$  & $e^{x + \ln y}$    & $1+3+1 = 5$   & 4 \\
\bottomrule
\end{tabular}
\end{center}
\end{definition}

\begin{remark}[Naive cost derivation for each pattern]
\label{rem:naive-derivations}
The naive costs in Definition~\ref{def:patterns} are computed by adding
SuperBEST v3 unit costs for each primitive in the expansion:
\begin{itemize}
  \item $\EML(x,y) = e^x - \ln y$: \texttt{exp} (1) + \texttt{ln} (1)
        + \texttt{sub} (2) $= 4$.
  \item $\EXL(x,y) = e^x \cdot \ln y$: \texttt{exp} (1) + \texttt{ln} (1)
        + \texttt{mul} (2) $= 4$.
  \item $\EAL(x,y) = e^x + \ln y$ (positive domain since $e^x > 0$):
        \texttt{exp} (1) + \texttt{ln} (1) + \texttt{add} (3) $= 5$.
  \item $\DEML(x,y) = e^{-x} - \ln y$: \texttt{neg} (2) + \texttt{exp} (1)
        + \texttt{ln} (1) + \texttt{sub} (2) $= 6$.
  \item $\LEAd(x,y) = \ln(e^x + y)$ (positive domain since $e^x > 0$):
        \texttt{exp} (1) + \texttt{add} (3) + \texttt{ln} (1) $= 5$.
  \item $\ELAd(x,y) = e^{x + \ln y} = e^x \cdot y$: \texttt{ln} (1)
        + \texttt{add} (3) + \texttt{exp} (1) $= 5$.
\end{itemize}
\end{remark}

\begin{definition}[Pattern Bonus]
\label{def:PB}
Let $M(E, \mathcal{P})$ be a \emph{maximal non-overlapping matching} of the
pattern set $\mathcal{P}$ against the expression~$E$: a set of sub-expression
sites $\{s_1, \ldots, s_m\}$ such that
\begin{enumerate}
  \item Each site $s_j$ is the root of a sub-expression of $E$ that matches
        some pattern $p_{i(j)} \in \mathcal{P}$.
  \item The sites are \emph{non-overlapping}: no node belongs to more than one
        matched sub-expression.
  \item The matching is \emph{maximal}: no additional site can be added without
        violating condition (2).
\end{enumerate}
For each matched site $s_j$, the per-site bonus is
$\bonus(p_{i(j)}) = k_{i(j)}^{\mathrm{naive}} - 1$.
The \emph{pattern bonus} is
\[
  \PB(E) \;=\; \sum_{j=1}^{m} \bonus(p_{i(j)})
         \;=\; \sum_{j=1}^{m} \bigl(k_{i(j)}^{\mathrm{naive}} - 1\bigr).
\]
\end{definition}

\begin{proposition}[$\PB(E) \geq 0$]
\label{prop:PB-nonneg}
For any expression $E$ and pattern set $\mathcal{P}$, $\PB(E) \geq 0$.
\end{proposition}

\begin{proof}
Each bonus term satisfies $k_i^{\mathrm{naive}} \geq 2 > 1$ for every pattern
in $\mathcal{P}$ (a pattern that maps to exactly one node must replace at least
two primitive nodes).  Hence every summand $\bonus(p_i) \geq 1 > 0$, and the
sum is non-negative.  If no patterns match, $m = 0$ and $\PB(E) = 0 \geq 0$.
\end{proof}

% ════════════════════════════════════════════════════════════════════════════
\section{Theorem T38 and Its Proof}
\label{sec:T38}
% ════════════════════════════════════════════════════════════════════════════

% ── 2.1 Statement ───────────────────────────────────────────────────────────
\subsection{Statement}

\begin{theorem}[T38 --- Cost Decomposition Theorem]
\label{thm:T38}
Let $E$ be any arithmetic expression computable by a finite EML operator tree
over $\mathcal{F}_{16}$.  Let $\mathcal{P}$ be the SuperBEST v3 compound
pattern set (Definition~\ref{def:patterns}) and let sharing be canonical
(share all identical sub-expressions).  Then:
\begin{equation}
  \boxed{\Cost(E) \;=\; \NC(E) \;-\; \SD(E) \;-\; \PB(E),}
\label{eq:T38}
\end{equation}
where $\NC(E) \geq 0$, $\SD(E) \geq 0$, and $\PB(E) \geq 0$.
\end{theorem}

% ── 2.2 Proof ────────────────────────────────────────────────────────────────
\subsection{Proof}

The proof proceeds in four steps.  We construct a sequence of expression
graphs, each at most as large as the previous, and track the exact node count
at each stage.

\begin{proof}
\textbf{Step 1 --- Na\"{i}ve tree and upper bound (T34).}

Construct the \emph{na\"{i}ve tree} $\Tnaive$ as follows.  For each occurrence
of a primitive operation $\mathrm{op}$ in the syntactic tree of~$E$, substitute
the operation with its SuperBEST v3 optimal sub-tree (1 node for \texttt{exp},
\texttt{ln}, \texttt{div}; 2 nodes for \texttt{neg}, \texttt{rec},
\texttt{mul}, \texttt{sub}; 3 nodes for \texttt{pow}, \texttt{add}).  Make no
attempt to share sub-expressions or to recognise compound patterns.  The result
$\Tnaive$ has exactly $\NC(E)$ nodes and correctly computes $E$.  Hence:
\[
  \Cost(E) \;\leq\; \NC(E).
\]

\textbf{Step 2 --- Sharing discount (construction and validity).}

Starting from $\Tnaive$, apply the following reductions exhaustively:
\begin{enumerate}
  \item \emph{Constant folding:} If a sub-tree has all leaf inputs equal to
        compile-time constants, replace it with a single constant leaf.
  \item \emph{Shared sub-expression elimination (CSE):} If a sub-tree $S$
        appears at $k \geq 2$ sites in the current graph, retain one copy and
        redirect all $k$ parent edges to that single copy.
\end{enumerate}
Each application of rule~(i) or~(ii) produces a valid DAG for $E$: rule~(i)
preserves the computed value by arithmetic correctness; rule~(ii) preserves
the computed value because both copies of $S$ produce identical outputs
(they are syntactically identical sub-trees over the same live-variable
inputs).

Let $G_1$ be the DAG resulting after applying rules~(i) and~(ii) to
exhaustion.  By Definition~\ref{def:SD}:
\[
  |G_1| \;=\; \NC(E) \;-\; \SD(E),
\]
where $|G_1|$ denotes the node count of $G_1$.  By
Proposition~\ref{prop:SD-nonneg}, $\SD(E) \geq 0$, so $|G_1| \leq \NC(E)$.

\textbf{Step 3 --- Pattern bonus (construction and validity).}

Starting from $G_1$, apply the compound-pattern substitutions from
$\mathcal{P}$ (Definition~\ref{def:patterns}) using a maximal non-overlapping
matching $M(E, \mathcal{P})$.  For each matched site $s_j$:
\begin{itemize}
  \item Remove the $k_{i(j)}^{\mathrm{naive}}$ primitive nodes at site $s_j$.
  \item Insert a single compound-operator node of type $p_{i(j)} \in
        \mathcal{F}_{16}$, wiring its inputs and output identically to the
        removed sub-tree.
\end{itemize}
Each such replacement is valid because $p_{i(j)}$ computes the same
function as the removed sub-tree (by definition of the compound pattern) and
$p_{i(j)} \in \mathcal{F}_{16}$ is a legal operator.

Let $G_2$ be the DAG after all pattern substitutions.  Then:
\[
  |G_2| \;=\; |G_1| \;-\; \PB(E) \;=\; \NC(E) \;-\; \SD(E) \;-\; \PB(E).
\]
By Proposition~\ref{prop:PB-nonneg}, $\PB(E) \geq 0$, so
$|G_2| \leq |G_1| \leq \NC(E)$.

\textbf{Step 4 --- Minimality and exhaustion of reductions.}

We claim $G_2 = \Gmin$, i.e.\ $\Cost(E) = |G_2|$.

First, $G_2$ is a valid DAG for $E$ (shown in Steps 2 and~3).

Second, $G_2$ is minimal.  Suppose for contradiction that there exists a valid
DAG $G'$ for $E$ with $|G'| < |G_2|$.  Then $G'$ must use either a further
sharing reduction or a further pattern substitution not present in $G_2$.
\begin{itemize}
  \item A further sharing reduction would mean $G_2$ contains a shared
        sub-expression that was not eliminated in Step~2.  But Step~2 was
        applied to exhaustion, so no such sub-expression exists.
  \item A further pattern substitution would mean $G_2$ contains a sub-tree
        matching some $p_i \in \mathcal{P}$ that was not covered by the
        maximal matching $M$.  But $M$ is maximal (Definition~\ref{def:PB}),
        so no uncovered matchable site exists.
  \item Any other reduction would constitute a new type of structural
        compression not captured by sharing or pattern recognition.  By the
        No Nesting Penalty lemma (established in the companion paper
        \emph{Cost\_Theory}), nesting depth introduces no overhead, so no
        interface nodes can be removed by reordering.  No other reduction type
        is structurally available in $\mathcal{F}_{16}$.
\end{itemize}
Hence no such $G'$ exists, and $G_2$ is minimal.  Therefore:
\[
  \Cost(E) \;=\; |G_2| \;=\; \NC(E) \;-\; \SD(E) \;-\; \PB(E).
\]

\textbf{Non-negativity.}
$\NC(E) \geq 0$ by definition (it is a sum of non-negative terms).
$\SD(E) \geq 0$ by Proposition~\ref{prop:SD-nonneg}.
$\PB(E) \geq 0$ by Proposition~\ref{prop:PB-nonneg}.
Since $\Cost(E) \geq 0$ (node counts are non-negative), the identity is
consistent.
\end{proof}

% ── 2.3 The three terms exhaust all reductions ───────────────────────────────
\subsection{Exhaustion of Reductions}

\begin{corollary}[No further reductions]
\label{cor:exhaustion}
After applying all canonical sharing reductions and all maximal non-overlapping
pattern substitutions, no node can be removed from the resulting DAG while
preserving the correct computation of~$E$.
\end{corollary}

\begin{proof}
Immediate from the minimality argument in Step~4 of the proof of
Theorem~\ref{thm:T38}.
\end{proof}

\begin{remark}[Gap between lower bounds and T38]
The lower bounds T35 ($\Cost(E) \geq 1$), T36 ($\Cost(E) \geq \lceil
d(E)/2\rceil$), and T37 ($\Cost(E) \geq n_{\exp} + n_{\ln}$) bracket
$\Cost(E)$ from below.  T38 pins the exact value.  The gap between any lower
bound $L$ and $\NC(E)$ is:
\[
  \NC(E) - L \;=\; \SD(E) + \PB(E) + \bigl(\Cost(E) - L\bigr).
\]
When $L = \Cost(E)$ exactly (the bound is tight), the gap equals
$\SD(E) + \PB(E)$.
\end{remark}

% ════════════════════════════════════════════════════════════════════════════
\section{Uniqueness of the Decomposition}
\label{sec:unique}
% ════════════════════════════════════════════════════════════════════════════

% ── 3.1 Canonical sharing strategy ──────────────────────────────────────────
\subsection{Canonical Sharing Strategy}

\begin{definition}[Canonical sharing strategy]
\label{def:canonical-sharing}
A sharing strategy is \emph{canonical} if it shares every pair of
syntactically identical sub-expressions that appear at two or more sites in the
expression.  Formally, two sub-expressions are identified if and only if they
have the same operator tree structure and the same leaf labels (variables or
constants).
\end{definition}

\begin{remark}[Uniqueness of canonical sharing]
Given a fixed expression $E$, the canonical sharing strategy produces a unique
DAG $G_{\mathrm{canon}}(E)$ (up to isomorphism): two implementations of the
canonical strategy must produce the same node-merging decisions because
structural equality of sub-expressions is a deterministic predicate.
\end{remark}

% ── 3.2 Uniqueness theorem ───────────────────────────────────────────────────
\subsection{Uniqueness Theorem}

\begin{theorem}[Uniqueness of T38 decomposition]
\label{thm:unique}
Fix a pattern set $\mathcal{P}$ and the canonical sharing strategy
(Definition~\ref{def:canonical-sharing}).  For any expression $E$:
\begin{enumerate}
  \item The value $\SD(E)$ is unique: every canonical sharing strategy produces
        the same discount.
  \item The value $\PB(E)$ is unique up to the choice of maximal non-overlapping
        matching.  If $\mathcal{P}$ contains no two patterns whose matches
        overlap in the expression, then $\PB(E)$ is fully determined.
  \item The identity \eqref{eq:T38} holds for every valid choice of maximal
        matching, and all such choices yield the same total $\PB(E)$.
\end{enumerate}
\end{theorem}

\begin{proof}
\textbf{(1) Uniqueness of $\SD(E)$.}
By Definition~\ref{def:SD}, $\SD(E) = \sum_{v: k_v \geq 2} (k_v - 1)\Cost(v)$.
The canonical sharing strategy identifies each sub-expression by its
syntactic structure.  The multiset of shared sub-expressions is thus uniquely
determined by~$E$.  The multiplicities $k_v$ are unique, and $\Cost(v)$ is
unique for each node $v$ (it is the optimal SuperBEST v3 cost of that
sub-expression).  Hence the sum is unique.

\textbf{(2) Uniqueness of $\PB(E)$.}
Suppose two maximal non-overlapping matchings $M$ and $M'$ of $\mathcal{P}$
against $E$ are given.  Both are maximal, so they cover the same total set of
pattern-matchable nodes (no site remains uncovered in either matching).

If patterns in $\mathcal{P}$ never overlap (i.e.\ no two patterns in
$\mathcal{P}$ can simultaneously match overlapping sub-expressions), then every
matchable node belongs to exactly one pattern site, and both $M$ and $M'$
assign the same site set.  The total bonus is therefore equal.

If patterns \emph{can} overlap (e.g.\ a node could be the root of both an
$\EML$ match and an $\EAL$ match), then different maximal matchings may assign
different bonuses.  In this case $\PB(E)$ depends on the matching, and one
conventionally takes the matching that maximises total bonus (greedy maximum
weight matching).

\textbf{(3) T38 under different matchings.}
Let $M$ and $M'$ be two maximal non-overlapping matchings with bonus totals
$B$ and $B'$.  Both $G_2$ (using $M$) and $G_2'$ (using $M'$) are valid DAGs
for $E$ with node counts $\NC(E) - \SD(E) - B$ and $\NC(E) - \SD(E) - B'$
respectively.  By the minimality argument of Theorem~\ref{thm:T38}, both are
minimal (they exhaust all reductions under the given matching).  If $B \neq
B'$, the two matchings produce DAGs of different sizes, contradicting
minimality of both.  Hence $B = B'$ and $\PB(E)$ is unique.
\end{proof}

% ── 3.3 Effect of extending the pattern set ──────────────────────────────────
\subsection{Effect of Extending the Pattern Set}

\begin{proposition}[Pattern set extension and decomposition shift]
\label{prop:extension}
Let $\mathcal{P}' \supset \mathcal{P}$ be an extended pattern set obtained by
adding patterns $q_1, \ldots, q_r$ to $\mathcal{P}$.  Let
$\PB'(E)$ be the pattern bonus under $\mathcal{P}'$.  Then:
\begin{enumerate}
  \item $\PB'(E) \;\geq\; \PB(E)$ for all $E$,
        with strict inequality when some new pattern $q_j$ matches in $E$.
  \item The decomposition remains valid:
        $\Cost(E) = \NC(E) - \SD(E) - \PB'(E)$,
        but $\SD(E)$ is unchanged (sharing is independent of pattern
        recognition).
  \item Uniqueness of $\PB'(E)$ holds under the same conditions as
        Theorem~\ref{thm:unique}.
\end{enumerate}
\end{proposition}

\begin{proof}
\textbf{(1).}
Any matching under $\mathcal{P}$ is also a valid (but not necessarily maximal)
matching under $\mathcal{P}'$.  The maximal matching under $\mathcal{P}'$ can
exploit additional patterns, so its total bonus is at least as large.

\textbf{(2).}
The construction in Theorem~\ref{thm:T38} (Steps~2 and~3) is valid for any
pattern set.  Step~2 (sharing) is independent of which patterns are used.
Step~3 uses whichever patterns are in $\mathcal{P}'$.  The minimality argument
in Step~4 applies unchanged, so the identity holds.

\textbf{(3).}
By the same argument as Theorem~\ref{thm:unique}(3), any two maximal matchings
under $\mathcal{P}'$ produce equal total bonuses, so $\PB'(E)$ is unique.
\end{proof}

\begin{remark}[Practical implication]
Adding new compound operators to $\mathcal{F}_{16}$ strictly increases
$\PB(E)$ for any expression containing the new pattern, and the decomposition
automatically accommodates the change without altering $\SD(E)$ or $\NC(E)$.
This modularity is a key structural property of T38.
\end{remark}

% ════════════════════════════════════════════════════════════════════════════
\section{Worked Examples}
\label{sec:examples}
% ════════════════════════════════════════════════════════════════════════════

We work three scientific examples in full detail, computing $\NC(E)$,
$\SD(E)$, $\PB(E)$, and $\Cost(E)$ and verifying the identity \eqref{eq:T38}.

% ── 4.1 Arrhenius rate law ────────────────────────────────────────────────────
\subsection{Example 1: Arrhenius Rate Law (Class A)}

\begin{example}[Arrhenius]
\label{ex:arrhenius}
The Arrhenius rate law is
\[
  k \;=\; A \cdot e^{-E_a / (RT)},
\]
where $A$ (pre-exponential factor), $E_a$ (activation energy), $R$ (gas
constant), and $T$ (temperature) are live variables or fixed parameters.

\textbf{Primitive decomposition.}
We treat $A$, $E_a$, $R$, $T$ as live variables (no constant folding).
The expression tree contains:
\[
  k \;=\; \texttt{mul}\bigl(A,\; \texttt{exp}(\texttt{neg}(\texttt{div}(E_a, \texttt{mul}(R,T))))\bigr).
\]
Counting operations: 1~\texttt{mul} (innermost), 1~\texttt{div},
1~\texttt{neg}, 1~\texttt{exp}, 1~\texttt{mul} (outer).

\textbf{Na\"{i}ve cost.}
\[
  \NC(k) \;=\; 2 + 1 + 2 + 1 + 2 \;=\; 8.
\]
(mul: 2, div: 1, neg: 2, exp: 1, mul: 2.)

\textbf{Sharing discount.}
No sub-expression appears more than once.  No constant folding applies (all
inputs live).  Hence $\SD(k) = 0$.

\textbf{Pattern bonus.}
Inspect the sub-expression $e^{-x}$ where $x = E_a/(RT)$.  This matches
$\DEML$ only if it is followed by $-\ln y$, which it is not here (there is no
$\ln$ in the expression).  No other pattern in $\mathcal{P}$ matches.
Hence $\PB(k) = 0$.

\textbf{Verification.}
\[
  \Cost(k) \;=\; \NC(k) - \SD(k) - \PB(k) \;=\; 8 - 0 - 0 \;=\; 8.
\]
The SuperBEST v3 optimal EML tree uses: $\LEdiv(E_a, \ELAd(R,T))$ for the
denominator ratio (2 nodes: 1 \texttt{div} via $\LEdiv$ + 2 for \texttt{mul}
via $\ELAd$), then \texttt{neg} (2 nodes), then \texttt{exp} (1 node), then
outer \texttt{mul} (2 nodes).  Total: $1 + 2 + 2 + 1 + 2 = 8$~nodes.
\hfill$\checkmark$
\end{example}

% ── 4.2 pH formula ────────────────────────────────────────────────────────────
\subsection{Example 2: pH Formula (Class C)}

\begin{example}[pH]
\label{ex:pH}
The pH formula is
\[
  \mathrm{pH} \;=\; -\log_{10}[\mathrm{H}^+]
              \;=\; -\frac{\ln[\mathrm{H}^+]}{\ln 10},
\]
where $[\mathrm{H}^+]$ is a live variable and $10$ is a numeric constant.

\textbf{Primitive decomposition.}
\[
  \mathrm{pH} \;=\; \texttt{neg}\!\left(\texttt{div}\!\left(\texttt{ln}([\mathrm{H}^+]),\; \texttt{ln}(10)\right)\right).
\]
Operations: 1~\texttt{ln} (argument), 1~\texttt{ln} (constant 10),
1~\texttt{div}, 1~\texttt{neg}.

\textbf{Na\"{i}ve cost.}
\[
  \NC(\mathrm{pH}) \;=\; 1 + 1 + 1 + 2 \;=\; 5.
\]
(ln: 1, ln: 1, div: 1, neg: 2.)

\textbf{Sharing discount.}
$\ln(10)$ is a compile-time constant (10 is fixed).  Constant folding replaces
$\ln(10) \approx 2.303$ with a single constant leaf, saving 1 node.
Hence $\SD(\mathrm{pH}) = 1$.

\textbf{Pattern bonus.}
After constant folding, the expression is
$\texttt{neg}(\texttt{div}(\texttt{ln}([\mathrm{H}^+]),\; c))$
where $c = \ln(10)$ is a constant.  The sub-expression
$\texttt{div}(\texttt{ln}(x), c)$ matches the pattern $\LEdiv(x, c)$ (a
single-node operator).  This is already cost~1 and counts as the \texttt{div}
(+\texttt{ln}) in the $\LEdiv$ compound: naive cost $= 1 + 1 = 2$, actual
cost $= 1$, bonus $= 1$.  However, $\LEdiv$ is the SuperBEST primitive for
division itself, so this saving is already captured in the unit cost
$c_{\mathrm{div}} = 1$.  No additional pattern from $\mathcal{P}$ matches.
Hence $\PB(\mathrm{pH}) = 0$.

\textbf{Verification.}
\[
  \Cost(\mathrm{pH}) \;=\; \NC(\mathrm{pH}) - \SD(\mathrm{pH}) - \PB(\mathrm{pH})
               \;=\; 5 - 1 - 0 \;=\; 4.
\]
The SuperBEST v3 tree uses: $\LEdiv([\mathrm{H}^+], c)$ (1 node) for
$\ln([\mathrm{H}^+])/\ln(10)$, then \texttt{neg} (2 nodes), for a total of
$1 + 2 = 3$ nodes over live variables, plus 1 constant-folded node for
$\ln(10)$.  Counting the materialized constant: $3 + 1 = 4$ nodes total.
\hfill$\checkmark$
\end{example}

% ── 4.3 Maxwell-Boltzmann speed distribution ─────────────────────────────────
\subsection{Example 3: Maxwell-Boltzmann Speed Distribution (Class D)}

\begin{example}[Maxwell-Boltzmann]
\label{ex:MB}
The Maxwell-Boltzmann probability density for molecular speeds is
\[
  f(v) \;=\; 4\pi \left(\frac{m}{2\pi k_B T}\right)^{3/2}
             v^2 \exp\!\left(-\frac{mv^2}{2k_BT}\right),
\]
where $v$ is the live variable (speed), $m$ is molecular mass, $k_B$ is the
Boltzmann constant, $T$ is temperature, and $4$, $\pi$, $2$ are numeric
constants.

\textbf{Simplification for analysis.}
Let $\alpha = m/(2k_BT)$ be pre-computed (constant folding of the constant
sub-expression $m/(2k_BT)$).  Then:
\[
  f(v) \;=\; C \cdot v^2 \cdot e^{-\alpha v^2},
\]
where $C = 4\pi(m/(2\pi k_BT))^{3/2}$ and $\alpha$ are fully determined by
constants.

\textbf{Primitive decomposition (after constant folding).}
The live-variable portion is $v^2 \cdot e^{-\alpha v^2}$:
\begin{align*}
  &\texttt{pow}(v, 2) \quad\text{(compute $v^2$)},\\
  &\texttt{mul}(\alpha, \texttt{pow}(v,2)) \quad\text{(compute $\alpha v^2$)},\\
  &\texttt{neg}(\texttt{mul}(\alpha, v^2)) \quad\text{(compute $-\alpha v^2$)},\\
  &\texttt{exp}(\texttt{neg}(\cdots)) \quad\text{(compute $e^{-\alpha v^2}$)},\\
  &\texttt{mul}(v^2, e^{-\alpha v^2}) \quad\text{(compute $v^2 e^{-\alpha v^2}$)},\\
  &\texttt{mul}(C, \cdots) \quad\text{(multiply by constant $C$)}.
\end{align*}

\textbf{Na\"{i}ve cost (live portion, $v^2$ counted twice).}
Without sharing: $2 \times \texttt{pow}(v,2)$ (each occurrence independent),
\texttt{mul} ($\alpha v^2$), \texttt{neg}, \texttt{exp}, \texttt{mul}
($v^2 \cdot e^{-\alpha v^2}$), \texttt{mul} ($C \cdot \cdots$):
\[
  \NC(f; \text{no sharing}) \;=\; 3 + 3 + 2 + 2 + 1 + 2 + 2 \;=\; 15.
\]
(Two \texttt{pow}: $3+3$; \texttt{mul}: 2; \texttt{neg}: 2; \texttt{exp}: 1;
\texttt{mul}: 2; \texttt{mul}: 2.)

\textbf{Sharing discount.}
$v^2$ appears twice (once as $\texttt{pow}(v,2)$ multiplied by $\alpha$, once
as the outer $v^2$ factor).  Sharing eliminates 1 copy of \texttt{pow}
(cost~3):
\[
  \SD(f) \;=\; (2-1) \cdot 3 \;=\; 3.
\]

\textbf{Pattern bonus.}
After sharing, inspect the sub-expression $e^{-\alpha v^2}$.  The pattern
$\DEML(x,y) = e^{-x} - \ln y$ does not match (no $\ln$ present).  The
pattern for $e^{-x}$ alone (i.e.\ \texttt{neg} followed by \texttt{exp}) is
not a compound node in $\mathcal{P}$ but rather a two-node chain with cost
$c_{\texttt{neg}} + c_{\texttt{exp}} = 2 + 1 = 3$ (No Nesting Penalty: no
overhead).  No $p_i \in \mathcal{P}$ matches any sub-expression here.
Hence $\PB(f) = 0$.

\textbf{Verification.}
\[
  \Cost(f) \;=\; \NC(f) - \SD(f) - \PB(f)
           \;=\; 15 - 3 - 0 \;=\; 12.
\]
The SuperBEST v3 tree: $v^2$ (1 shared \texttt{pow} node = 3), then
$\alpha v^2$ (1 \texttt{mul} = 2), then $-\alpha v^2$ (\texttt{neg} = 2),
then $e^{-\alpha v^2}$ (\texttt{exp} = 1), then $v^2 \cdot e^{-\alpha v^2}$
(\texttt{mul} = 2), then $C \cdot \cdots$ (\texttt{mul} = 2).
Total over live variables: $3 + 2 + 2 + 1 + 2 + 2 = 12$~nodes.
\hfill$\checkmark$
\end{example}

% ── 4.4 Summary table ────────────────────────────────────────────────────────
\subsection{Summary of Worked Examples}

\begin{table}[h]
\centering
\caption{T38 decomposition for the three worked examples.}
\label{tab:examples}
\smallskip
\begin{tabular}{@{}lccccl@{}}
\toprule
\textbf{Equation} & \textbf{Class} & $\NC$ & $\SD$ & $\PB$ & $\Cost$ \\
\midrule
Arrhenius $k = Ae^{-E_a/RT}$       & A & 8  & 0 & 0 & 8  \\
pH $= -\ln[\mathrm{H}^+]/\ln 10$   & C & 5  & 1 & 0 & 4  \\
Maxwell-Boltzmann $f(v)$            & D & 15 & 3 & 0 & 12 \\
\bottomrule
\end{tabular}
\end{table}

\begin{remark}[When $\PB > 0$]
None of the three examples activates a pattern bonus because none contains a
matching compound sub-expression from $\mathcal{P}$.  An example with
$\PB > 0$ would be the softmax logit
$\sigma_i = e^{x_i} / \sum_j e^{x_j}$, where $\DEML$ patterns appear in the
denominator terms $e^{x_j - x_i}$ after the log-sum-exp trick, yielding
$\PB > 0$ for the full softmax probability.
\end{remark}

% ════════════════════════════════════════════════════════════════════════════
\section*{Acknowledgments}
% ════════════════════════════════════════════════════════════════════════════

This document is the formal companion to Section~2 of \emph{The General Theory
of SuperBEST Cost} (T34--T39, April 2026).  It provides the self-contained
proof record for Theorem T38.

\bigskip
\noindent\textit{Almaguer, A.R.\ (2026).
T38: The Cost Decomposition Theorem.
monogate research. \url{https://monogate.org}}

\end{document}
